%Página
\section{Autoridade de Certificação (CA)}

\subsection{O que é}

Uma Autoridade de Certificação (CA, do inglês Certificate Authority) é uma entidade confiável que emite certificados digitais. Esses certificados são usados para verificar a identidade de uma entidade (pessoa, organização ou dispositivo) e para estabelecer a confiança nas comunicações digitais. A CA desempenha um papel crucial na infraestrutura de chave pública (PKI, do inglês Public Key Infrastructure), que é usada para proteger transações e comunicações na internet.

\subsection{Funções Principais de uma CA}

\begin{itemize}
	\item \textbf{Emissão de Certificados Digitais: }A principal função de uma CA é emitir certificados digitais. 
	\item \textbf{Verificação de Identidade: }Antes de emitir um certificado, a CA verifica a identidade da entidade que solicita o certificado.
	\item \textbf{Assinatura de Certificados: }A CA assina digitalmente o certificado usando sua chave privada, garantindo que o certificado não tenha sido adulterado e que a identidade foi verificada pela CA.
	\item \textbf{Gestão de Certificados: }A CA também gerencia o ciclo de vida dos certificados.
\end{itemize}

\pagebreak
